% !TeX program = xelatex
\documentclass[10pt]{article}

\usepackage{hyperref}
\usepackage{xcolor}
\usepackage{calc}
\usepackage{graphicx}
\usepackage{tikz}
\usepackage{fontspec}
\usepackage{fontawesome5}
\usepackage{titlesec}
\usepackage{enumitem}
\usepackage{fancybox}
\usepackage{xeCJK}

\hypersetup{hidelinks}
\urlstyle{same}

%%%%%%%%%%%%%%%%%%%%
% 设置
%%%%%%%%%%%%%%%%%%%%

\providecommand{\ResumeItemSep}{0.12em}
\providecommand{\ResumeTopSep}{0.18em}
\providecommand{\ResumeArrayStretch}{1.12}
\providecommand{\ResumeLineSpread}{1.08}
\providecommand{\ResumeSectionBefore}{0.55em}
\providecommand{\ResumeSectionAfter}{0.28em}

\setlength{\parindent}{0pt}
\setlength{\parskip}{0pt}
\pagenumbering{gobble}
\setlist[itemize]{
    leftmargin=1.35em,
    itemsep=\ResumeItemSep,
    topsep=\ResumeTopSep,
    parsep=0pt,
    partopsep=0pt
}
\setlist[enumerate]{leftmargin=*}
\renewcommand{\arraystretch}{\ResumeArrayStretch}
\linespread{\ResumeLineSpread}

\titleformat{\section}
  {\Large\bfseries\raggedright}
  {}{0em}
  {}
  [{\color{secondarycolor}\titlerule}]
\titlespacing*{\section}{0pt}{\ResumeSectionBefore}{\ResumeSectionAfter}

\usepackage[
    a4paper,
    left=1.1cm,
    right=1.1cm,
    top=0.75cm,
    bottom=0.85cm,
    nohead
]{geometry}

% 字体设置
% 说明：Noto Serif SC 没有原生斜体字型。将 CJK 的 italic 字型映射到常规字形
%（与过去“找不到斜体、默认替换为常规字形”的实际渲染一致），可消除
% “Font shape `TU/NotoSerifSC(0)/m/it' undefined” 告警；如需真正斜体观感，
% 可在此行追加 ItalicFeatures={FakeSlant=0.18}。
\setCJKmainfont[
    Path=fonts/,
    Extension=.otf,
    BoldFont=*-Bold,
    ItalicFont={NotoSerifSC},
]{NotoSerifSC}

% 自定义颜色
\definecolor{primarycolor}{RGB}{200, 60, 60}
\definecolor{secondarycolor}{RGB}{200, 68, 28}

\newlength{\iconwidth}
\setlength{\iconwidth}{1.5em}

\newcommand{\sectionicon}[2]{\makebox[\iconwidth][c]{\color{primarycolor}{#1}}\quad #2}

\newcommand{\projectheading}[3]{%
    \noindent\makebox[\textwidth][s]{%
        \makebox[0pt][l]{{\large \textbf{#1}}}%
        \hfill
        \makebox[0pt][c]{{\normalsize \textbf{#2}}}%
        \hfill
        \makebox[0pt][r]{#3}%
    }%
}

% 通用个人信息：修改这里即可同步到所有岗位模板。
\newcommand{\ResumeName}{无绝}
\newcommand{\ResumeBirthDate}{2002.10}
\newcommand{\ResumeHometown}{湖北咸宁}
\newcommand{\ResumeSchool}{重庆邮电大学}
\newcommand{\ResumeEnglishLevel}{CET-6}
\newcommand{\ResumeEmail}{250xxxxxxx@qq.com}
\newcommand{\ResumePhone}{198xxxxxxxx}
\newcommand{\ResumeHomepageLabel}{github.com/Dithob}
\newcommand{\ResumeHomepageUrl}{https://github.com/Dithob}
\newcommand{\ResumePhoto}{images/xjpic.jpg}
\newcommand{\ResumeHeaderImage}{images/foot.png}

% 通用教育背景：修改这里即可同步到所有岗位模板。
\newcommand{\ResumeGraduateSchool}{重庆邮电大学}
\newcommand{\ResumeGraduateMajor}{计算机技术}
\newcommand{\ResumeGraduateDegree}{硕士}
\newcommand{\ResumeGraduateDate}{2024.09--2027.06}
\newcommand{\ResumeGraduateAwards}{特等学业奖学金 $\times$ 1、二等学业奖学金 $\times$ 1}

\newcommand{\ResumeUndergraduateSchool}{重庆邮电大学}
\newcommand{\ResumeUndergraduateMajor}{计算机科学与技术}
\newcommand{\ResumeUndergraduateDegree}{本科}
\newcommand{\ResumeUndergraduateDate}{2020.09--2024.07}
\newcommand{\ResumeUndergraduateAwards}{三等学业奖学金；GPA: 3.28/4.0（专业前 25\%）}

\newcommand{\ResumeEducationBlock}{%
    {\textbf{\ResumeGraduateSchool}}，\ResumeGraduateMajor，\textit{\ResumeGraduateDegree} \hfill \ResumeGraduateDate
    \begin{itemize}
        \item \textbf{荣誉奖项}：\ResumeGraduateAwards
    \end{itemize}

    \vspace{0.2em}
    {\textbf{\ResumeUndergraduateSchool}}，\ResumeUndergraduateMajor，\textit{\ResumeUndergraduateDegree} \hfill \ResumeUndergraduateDate
    \begin{itemize}
        \item \textbf{荣誉奖项}：\ResumeUndergraduateAwards
    \end{itemize}
}


%%%%%%%%%%%%%%%%%%%%
% 文章内容
%%%%%%%%%%%%%%%%%%%%

\newcommand{\ResumeTargetRole}{测试开发}
\providecommand{\ResumeTargetRole}{求职方向待定}

\newcommand{\ResumeContact}{%
    \footnotesize
    \textcolor{white}{%
        \faEnvelope \quad 邮箱：\href{mailto:\ResumeEmail}{\ResumeEmail}
        \hspace{3.2em}
        \faPhone \quad 电话：\ResumePhone
        \hspace{3.2em}
        \faGithub \quad 主页：\href{\ResumeHomepageUrl}{\ResumeHomepageLabel}
    }%
}

\newcommand{\ResumeContactBar}{%
    \begin{tikzpicture}[remember picture, overlay]
        \node[anchor=north, inner sep=0pt](headercontacts) at (current page.north){%
            \includegraphics[width=\paperwidth]{\ResumeHeaderImage}
        };
        \node[anchor=center] at(headercontacts.center){\ResumeContact};
    \end{tikzpicture}
    \vspace{-1.15em}
}

\newcommand{\ResumeProfileSection}{%
    \section[个人信息]{\sectionicon{\faAddressCard}{个人信息}}
    \begin{tabular*}{\textwidth}{@{\extracolsep{\fill}}lll}
        \textbf{姓\qquad 名}：\ResumeName & \textbf{出生年月}：\ResumeBirthDate & \textbf{籍\qquad 贯}：\ResumeHometown \\
        \textbf{学\qquad 校}：\ResumeSchool & \textbf{英语水平}：\ResumeEnglishLevel & \textbf{求职方向}：\ResumeTargetRole
    \end{tabular*}
}

\newcommand{\ResumeEducationSection}{%
    \section[教育背景]{\sectionicon{\faGraduationCap}{教育背景}}
    \ResumeEducationBlock
}

\newcommand{\ResumeProfileAndEducation}{%
    %%%%%%%%%%%%%%%%%%%%
    % 个人信息与教育背景
    %%%%%%%%%%%%%%%%%%%%

    \begin{minipage}[t]{0.79\textwidth}
        \ResumeProfileSection

        \ResumeEducationSection
    \end{minipage}
    \hfill
    \begin{minipage}[t]{0.14\textwidth}
        \vspace{-0.35em}
        \setlength{\fboxsep}{0pt}
        % 双线框 (\doublebox) 会给照片额外加约 4pt 边框宽；图片若占满 \linewidth 会触发
        % “Overfull \hbox (3.99988pt too wide)” 告警，故图片宽度预留 4pt，使含框总宽恰好等于栏宽。
        \doublebox{\includegraphics[width=\dimexpr\linewidth-4pt\relax]{\ResumePhoto}}
    \end{minipage}
}


\begin{document}

    \ResumeContactBar
    \ResumeProfileAndEducation
    % \vspace{0.45em}

    %%%%%%%%%%%%%%%%%%%%
    % 专业技能
    %%%%%%%%%%%%%%%%%%%%

    \section[专业技能]{\makebox[\iconwidth][c]{\color{primarycolor}{\faWrench}}\quad 专业技能}
    \begin{itemize}
        \item \textbf{测试开发基础}：熟练运用等价类、边界值、场景法、状态机、决策表等测试设计方法；熟悉需求评审、用例评审、提测执行、缺陷跟进、回归验证等测试全流程；熟练使用 TAPD、智研平台管理用例与缺陷。
        \item \textbf{工程与平台能力}：熟练使用 Python、Pytest、Requests 编写接口自动化，Playwright、Selenium、uiautomator2、ADB 编写 UI/移动端自动化，JMeter、Locust 编写性能压测；熟悉 Linux、FastAPI、MySQL、Docker、Git 等工程能力。
        \item \textbf{AI 应用与测试能力}：理解 RAG、Harness、Memory 等 Agent 机制，熟悉 MCP、Skill、Subagent 等 Agent 工具框架，了解 Skill 等 AI 工具的能力边界及测试方法；能构建 Golden Dataset 与 Bad Case 回流闭环，设计幻觉、拒答、格式化、多轮回归等测试集，以关键词、JSON Schema、相似度、规则断言批量校验；有 Agent 用例自动更新与人工复核的落地经验，日常以 AI 编码助手完成测试脚本与代码审查。
    \end{itemize}

    % \vspace{0.35em}

    %%%%%%%%%%%%%%%%%%%%
    % 实习经历
    %%%%%%%%%%%%%%%%%%%%

    \section[实习经历]{\makebox[\iconwidth][c]{\color{primarycolor}{\faBriefcase}}\quad 实习经历}

    \projectheading{腾讯云与智慧产业研发公司}{测开实习生}{2026.06--2026.09}
    \begin{itemize}
        \item \textbf{手机地图 UI 自动化测试 Agent}：参与腾讯地图 UI 自动化测试平台建设，负责 Agent 编排、Skill 与 MCP 工具落地，覆盖 17 个业务域、600+ 用例。
        \item \textbf{多 Agent 交叉检查与 Skill 评测}：以多智能体独立生成、交叉比对合并产出 85 条分级用例，结合人工复核形成“AI 生成 + 人工判断”闭环；为地图数据查询、地图渲染两类 Skill 规划 120 条评测用例，沉淀 Agent/Skill 评测方法论。
        \item \textbf{腾讯地图发版需求测试}：完成【你用我赔 H5】发版需求测试分析，识别跨天状态机、重复进入、端外分享、埋点与兼容性风险，建立 R1--R8 风险清单并与产品/研发逐项对齐；按智研平台规范编写用例、为研发创建自测计划。
    \end{itemize}

    \vspace{0.25em}

    % \projectheading{重庆啄木鸟网络科技有限公司}{算法实习生（AI 应用测试与交付）}{2025.06--2025.09}
    % \begin{itemize}
    %     \item \textbf{核心职责}：负责大模型应用研发与交付验证，覆盖模型服务接口联调、Prompt/RAG 流程测试、功能冒烟与回归用例设计、接口异常排查和线上 Bad Case 分析，推动模型回复质量与服务稳定性优化。
    %     \item \textbf{小啡 GPT 智能回复模型}：围绕家电客服多轮对话构建意图识别、槽位抽取、知识命中、指令遵循和结构化输出测试集；使用脚本批量调用模型 API，对回复做关键词、JSON Schema、相似度和规则断言校验，沉淀 100+ 条对话回归样例，定位上下文丢失、知识召回错误和格式漂移等问题。
    %     \item \textbf{智能家电识别与自动入库系统}：针对 OCR/视觉大模型识别、联网补全、向量去重和入库链路设计接口测试与数据质量校验；构造正常、缺字段、重复商品、低置信度识别等场景，通过数据库断言验证字段完整性、去重结果和入库状态，提升交付前缺陷发现效率。
    % \end{itemize}

    % \vspace{0.35em}

    %%%%%%%%%%%%%%%%%%%%
    % 项目经历
    %%%%%%%%%%%%%%%%%%%%

    \section[项目经历]{\makebox[\iconwidth][c]{\color{primarycolor}{\faProjectDiagram}}\quad 项目经历}

    \projectheading{智能持续测试平台}{独立设计与开发}{2025.10--2026.02}
    \begin{itemize}
        \item \textbf{平台与框架设计}：基于 Python/Pytest/Requests 搭建五层接口自动化框架，基于 OpenAPI 自动生成契约用例校验状态码、响应 Schema 与业务错误码；以 Vue.js + FastAPI + MySQL + Redis + Docker 构建轻量测试平台，统一管理项目与用例。
        \item \textbf{UI 自动化与性能压测}：基于 Playwright 落地 PO 模式业务用例，集成 Trace、截图与录制支持回放；使用 JMeter/Locust 设计阶梯加压、稳定性与峰值压测，统计 QPS、P95、错误率并定位慢接口。
        \item \textbf{AI 增强与质量闭环}：将 Pytest/Playwright/Locust 执行封装为平台任务，经 Redis 队列异步调度；引入 LLM 用例生成流水线（生成草稿、人工复核、留存 Bad Case）；报告页聚合通过率、失败原因与趋势，形成测试闭环。
    \end{itemize}

    \vspace{0.25em}

    \projectheading{腾讯地图 UI 自动化测试 Agent}{参与框架建设}{2026.06--2026.09}
    \begin{itemize}
        \item \textbf{智能体编排与用例生成}：构建智能体测试体系，编排 13 个 Skill 与 3 个自研 MCP 服务，打通需求理解到精准回归的全流程，覆盖 17 个业务域、600+ 用例；联动 TAPD/Figma 自动生成可执行 case.py，单用例生成约 10--15 分钟（手写需 1--2 小时）。
        \item \textbf{MCP 工具开发与多模态感知}：自研设备控制与失败诊断 MCP，接入地图 MCP 解析地点坐标；融合 UIAutomator、OCR、OmniParser、OpenCV 四层感知，链式降级解决自绘控件无 resource-id 的定位问题；scrcpy 投屏将截图耗时从 1s 降至 30--100ms。
        \item \textbf{知识工程与语义检索}：设计纯文件知识图谱（Page/Component/Pitfall/Operation 四类节点），基于 23 个研发仓库 UI 索引做静态推理，让 Agent 写用例前先读懂研发代码；以 AST 扫描 1500+ 私有 helper，结合 jieba + BM25 实现公共方法语义检索与经验复用。
        \item \textbf{可靠性工程与自愈闭环}：以意图锁定、九维度自审与人工接管等护栏保障生成质量；Lint 门禁与逐步断言保证代码稳定，health 回测通过率淘汰 flaky 用例；失败后基于证据包自动再生成补丁形成自愈闭环。
    \end{itemize}

    % \vspace{0.35em}

    %%%%%%%%%%%%%%%%%%%%
    % 在校成果
    %%%%%%%%%%%%%%%%%%%%

    \section[在校成果]{\makebox[\iconwidth][c]{\color{primarycolor}{\faAward}}\quad 在校成果}
    \begin{itemize}
        \item \textbf{竞赛}：2025 “互联网+”创新创业大赛银奖；中国国际大学生创新大赛校级二等奖；重庆市 AI 大模型创新应用大赛省级三等奖
        \item \textbf{论文}：《LHAF: A Lightweight Hierarchical Adaptive Framework for LLM-based Multimodal Emotion Recognition in Conversations》（SCI I 区在投）
    \end{itemize}

\end{document}
